\pnum{5.1}{Two-particle separation into center-of-mass and relative motion}
Put $M=m_1+m_2$, $\mu=m_1m_2/M$,
$\mathbf r=\mathbf r_1-\mathbf r_2$ and
$\mathbf R=(m_1\mathbf r_1+m_2\mathbf r_2)/M$.
Solving these two linear vector equations gives
\begin{equation}
\boxed{\mathbf r_1=\mathbf R+\frac{\mu}{m_1}\mathbf r,\qquad
\mathbf r_2=\mathbf R-\frac{\mu}{m_2}\mathbf r}.
\end{equation}
The multivariable chain rule gives
$\nabla_1=(m_1/M)\nabla_R+\nabla_r
=(\mu/m_2)\nabla_R+\nabla_r$, and
$\nabla_2=(\mu/m_1)\nabla_R-\nabla_r$.
Substitute into $-\sh^2(\nabla_1^2/m_1+
\nabla_2^2/m_2)/2$. The mixed derivatives cancel:
\begin{equation}
\boxed{H=-\frac{\sh^2}{2M}\nabla_R^2
-\frac{\sh^2}{2\mu}\nabla_r^2+V(\mathbf r)}.
\end{equation}
For $\Psi(\mathbf R,\mathbf r)=F(\mathbf R)u(\mathbf r)$,
division by $Fu$ separates constants:
\begin{align}
-\frac{\sh^2}{2M}\nabla_R^2F&=E_{\rm CM}F,\\
\left[-\frac{\sh^2}{2\mu}\nabla_r^2+V(\mathbf r)\right]u&=E_{\rm rel}u,
\end{align}
where $\boxed{E=E_{\rm CM}+E_{\rm rel}}$.
This exactly maps a central two-body problem
onto a free center of mass and a single reduced-mass particle.

\pnum{5.2}{Reduced mass: isotope shifts, positronium and muonic hydrogen}
For an orbiting particle of mass $m$ and nuclear mass $M$,
$\mu=mM/(m+M)$. Coulomb energies scale with $\mu$,
so using $m$ in place of $\mu$ overestimates
the binding energy fractionally by
\begin{equation}
\frac{m-\mu}{\mu}=\frac mM.
\end{equation}
For ordinary hydrogen $m_e/m_p\simeq1/1836$,
giving $\boxed{0.054\%}$ (two significant figures).
For a Balmer transition $3\to2$,
\begin{equation}
\frac1\lambda=
R_\infty\frac{\mu}{m_e}
\left(\frac14-\frac19\right)
=\frac{5R_\infty\mu}{36m_e},
\quad
\lambda(M)=\frac{36m_e}{5R_\infty\mu(M)}.
\end{equation}
Because deuterium's nucleus is heavier,
its $\mu$ is larger, so its red line has a shorter
wavelength. Inserting $M_D\simeq2m_p$ gives
$\boxed{\lambda_H-\lambda_D\simeq0.18\ {\rm nm}}$.
For positronium, $\mu=m_e/2$, so
$\boxed{|E_1|=6.80\ {\rm eV}}$.
For muonic hydrogen $\mu_\mu=
m_\mu m_p/(m_\mu+m_p)$, with
$m_\mu=206.77m_e$, the reduced-mass ratio
$\mu_\mu/\mu_H\simeq186$;
the Lyman-$\alpha$ wavelength scales inversely
with $\mu$, giving
$\boxed{\lambda_\mu\simeq0.653\ {\rm nm}}$.
\checkit{The last value uses the nonrelativistic
point-nucleus spectrum; finite proton size and
QED corrections matter for precision muonic spectroscopy.}

\pnum{5.3}{Isotope splitting in the vibrational spectrum of HCl}
A diatomic oscillator has $\omega=\sqrt{k/\mu}$,
$\mu=m_Hm_{\rm Cl}/(m_H+m_{\rm Cl})$.
For approximate isotope masses $m_H=1$,
$m_{35}=35$, and $m_{37}=37$ (atomic units),
\begin{equation}
\mu_{35}=\frac{35}{36},\qquad
\mu_{37}=\frac{37}{38}.
\end{equation}
The frequencies are
$\nu_{35}/\nu_{37}=\sqrt{\mu_{37}/\mu_{35}}
=\sqrt{1332/1330}=1.0007516\ldots$,
so
\begin{equation}
\boxed{\Delta\nu\simeq7.51\times10^{-4}\nu}.
\end{equation}
The vibrational line for the heavier chlorine isotope
has the slightly lower frequency.

\pnum{5.4}{Normalization of bosonic and fermionic two-orbital states}
Take normalized orbitals $a,b$, with
$S=\bracket a b$. For
$\Psi_\pm=A[a(1)b(2)\pm b(1)a(2)]$,
\begin{equation}
\|\Psi_\pm\|^2=
|A|^2[2\pm2|S|^2].
\end{equation}
For orthogonal orbitals $S=0$,
$\boxed{A=1/\sqrt2}$ for both signs.
For the same normalized orbital $a=b$,
the symmetric bracket is $2a(1)a(2)$:
$\boxed{A=1/2}$.
The antisymmetric bracket is identically zero,
so no normalized fermion state exists.

\pnum{5.5}{Excited levels of two particles in a square well}
Let $\epsilon=\pi^2\sh^2/(2ma^2)$,
$\phi_n(x)=\sqrt{2/a}\sin(n\pi x/a)$.
For particles on $0<x_1,x_2<a$,
\begin{equation}
H=-\frac{\sh^2}{2m}
(\partial_{x_1}^2+\partial_{x_2}^2),
\quad E_{nm}=(n^2+m^2)\epsilon.
\end{equation}
For $n\ne m$, define
$\Phi^\pm_{nm}=
[\phi_n(1)\phi_m(2)\pm\phi_m(1)\phi_n(2)]/\sqrt2$.
Because the two product terms have the same sum
of one-body energies,
$H\Phi^-_{12}=5\epsilon\Phi^-_{12}$:
the fermionic ground state is verified.
Beyond the levels explicitly displayed in the
textbook example, the next two \emph{distinct}
levels are:
\begin{center}\begin{tabular}{lccc}\toprule
Statistics&Level&Wavefunction(s)&Degeneracy\\\midrule
Distinguishable&$8\epsilon$&$\phi_2(1)\phi_2(2)$&1\\
&$10\epsilon$&$\phi_1(1)\phi_3(2)$,
$\phi_3(1)\phi_1(2)$&2\\
Bosons&$8\epsilon$&$\phi_2(1)\phi_2(2)$&1\\
&$10\epsilon$&$\Phi^+_{13}$&1\\
Fermions&$10\epsilon$&$\Phi^-_{13}$&1\\
&$13\epsilon$&$\Phi^-_{23}$&1\\\bottomrule
\end{tabular}\end{center}
For fermions, $(2,2)$ is forbidden in this
spin-ignored model. Different ordered products
count as different states only when the
particles are distinguishable.

\pnum{5.6}{Exchange correction to mean squared separation in a well}
For $n\ne l$, the one-body moments are
\begin{align}
\av{x}_n&=a/2,&
\av{x^2}_n&=a^2\left(\frac13-\frac1{2\pi^2n^2}\right),\\
X_{nl}&=\int_0^a x\,\phi_n(x)\phi_l(x)\dd x\\
&=\frac{a}{\pi^2}\left[
\frac{(-1)^{n-l}-1}{(n-l)^2}
-\frac{(-1)^{n+l}-1}{(n+l)^2}\right].
\end{align}
Integrate $x\cos(k\pi x/a)$ by parts to
derive the last line; $X_{nl}=0$ for equal
parity, otherwise
$X_{nl}=-8anl/[\pi^2(n^2-l^2)^2]$.
For distinguishable particles,
$\av{(x_1-x_2)^2}_{D}=
\av{x^2}_n+\av{x^2}_l-2\av{x}_n\av{x}_l$:
\begin{equation}
\boxed{D_{nl}=a^2\left[
\frac16-\frac1{2\pi^2}
\left(\frac1{n^2}+\frac1{l^2}\right)\right]}.
\end{equation}
The exchange cross term in
$\av{x_1x_2}$ is $\pm|X_{nl}|^2$,
so
\begin{equation}
\boxed{\av{(x_1-x_2)^2}_{B}=D_{nl}-2|X_{nl}|^2,\quad
\av{(x_1-x_2)^2}_{F}=D_{nl}+2|X_{nl}|^2}.
\end{equation}
Bosons are closer on average in this
specific observable, and fermions farther apart.
This arises from exchange symmetry,
not an additional mechanical force.

\pnum{5.7}{Two oscillator orbitals: position densities and relative separation}
Let $b=\sqrt{\sh/(m\omega)}$,
$\phi_0(x)=\pi^{-1/4}b^{-1/2}
\exp[-x^2/(2b^2)]$,
and $\phi_1(x)=\sqrt2(x/b)\phi_0(x)$.
With one particle in each orbital,
\begin{equation}
\Psi_D=\phi_0(x_1)\phi_1(x_2),\qquad
\Psi_\pm=\frac{\phi_0(x_1)\phi_1(x_2)
\pm\phi_1(x_1)\phi_0(x_2)}{\sqrt2}.
\end{equation}
The distinguishable, boson, and fermion
probability surfaces are $|\Psi_D|^2$,
$|\Psi_+|^2$ and $|\Psi_-|^2$;
the last vanishes at $x_1=x_2$.
Since $\av{x^2}_0=b^2/2$,
$\av{x^2}_1=3b^2/2$ and
$\av{x}_{01}=b/\sqrt2$,
\begin{equation}
\boxed{\av{(x_1-x_2)^2}_D=2b^2,\quad
\av{(x_1-x_2)^2}_B=b^2,\quad
\av{(x_1-x_2)^2}_F=3b^2}.
\end{equation}
Introduce $R=(x_1+x_2)/2$ and
$r=x_1-x_2$, so
$x_1=R+r/2$, $x_2=R-r/2$.
The squared amplitudes factor into
a common Gaussian
$\exp[-2R^2/b^2-r^2/(2b^2)]$ times
$(R-r/2)^2$ for $D$, $4R^2$ for $B$,
or $r^2$ for $F$ (up to constants).
After integrating over $R$, put
$g_b(r)=\exp[-r^2/(2b^2)]/(\sqrt{2\pi}b)$.
The signed relative-displacement densities are
\begin{equation}
\boxed{p_D(r)=\tfrac12(1+r^2/b^2)g_b(r),\quad
p_B(r)=g_b(r),\quad
p_F(r)=(r^2/b^2)g_b(r)}.
\end{equation}
The distribution of distance $d=|r|\ge0$
is $\boxed{P_j(d)=2p_j(d)}$.
All three integrate to one, and the
second moments recover the results above.
Finally, the one-body number-density operator
$n(x)=\sum_{j=1}^2\delta(x-x_j)$ gives
the \emph{same} mean in every case:
\begin{equation}
\boxed{\av{n(x)}=
|\phi_0(x)|^2+|\phi_1(x)|^2}.
\end{equation}
Thus one-body density cannot distinguish
exchange symmetry here; two-body correlations can.

\pnum{5.8}{Three-particle bosonic permanent and fermionic determinant}
For three orthonormal orbitals $a,b,c$,
a particular distinguishable assignment is
$\boxed{\Psi_D=a(1)b(2)c(3)}$.
The completely symmetric state is
\begin{equation}
\boxed{\Psi_B=\frac1{\sqrt{3!}}
\sum_{\pi\in S_3}
a(\pi(1))b(\pi(2))c(\pi(3))}.
\end{equation}
Every one of the $3!$ terms is orthogonal
to all others, so the norm is $3!/3!=1$;
permutations merely reorder the sum.
The completely antisymmetric state is
\begin{equation}
\boxed{\Psi_F=\frac1{\sqrt{3!}}
\begin{vmatrix}
a(1)&b(1)&c(1)\\
a(2)&b(2)&c(2)\\
a(3)&b(3)&c(3)
\end{vmatrix}}.
\end{equation}
Exchanging two particle labels swaps
two rows and reverses the sign; repeated
orbitals give identical columns and zero.
